\documentclass[12pt]{article}

\usepackage[T1]{fontenc}
\usepackage[utf8]{inputenc}
\usepackage[romanian]{babel}

\usepackage{amsmath, amssymb, amsthm}
\usepackage{geometry}
\geometry{a4paper, margin=2.5cm}

\title{Variabile Aleatoare și Vectori Aleatori}
\author{}
\date{}

\theoremstyle{plain}
\newtheorem{theorem}{Teoremă}
\newtheorem{proposition}{Propoziție}
\newtheorem{corollary}{Corolar}
\newtheorem{definition}{Definiție}

\begin{document}

\maketitle

\section{Variabile aleatoare}

Fie $(\Omega, \mathcal{F}, \mathbb{P})$ un spațiu de probabilitate.
O aplicație



\[
X : \Omega \to \mathbb{R}
\]



se numește \textbf{variabilă aleatoare} dacă este măsurabilă, adică:



\[
X^{-1}(B) \in \mathcal{F}
\quad \text{pentru orice } B \in \mathcal{B}(\mathbb{R}).
\]



Aceasta înseamnă că preimaginile mulțimilor boreliene sunt evenimente.

\section{Criteriu de măsurabilitate}

\begin{proposition}
O aplicație $X : \Omega \to \mathbb{R}$ este variabilă aleatoare dacă și numai dacă
pentru orice mulțime de forma



\[
(-\infty, a], \qquad a \in \mathbb{R},
\]



avem:



\[
X^{-1}((-\infty, a]) \in \mathcal{F}.
\]



Mulțimile de forma $(-\infty, a]$ generează sigma-algebra boreliană
$\mathcal{B}(\mathbb{R})$, deci este suficient să verificăm măsurabilitatea
doar pe această familie generatoare.
\end{proposition}

\section{Vectori aleatori}

Pentru $n \ge 2$, o aplicație



\[
X = (X_1, X_2, \dots, X_n) : \Omega \to \mathbb{R}^n
\]



se numește \textbf{vector aleator} dacă este măsurabilă față de
$\mathcal{F}$ și $\mathcal{B}(\mathbb{R}^n)$, adică:



\[
X^{-1}(B) \in \mathcal{F}
\quad \text{pentru orice } B \in \mathcal{B}(\mathbb{R}^n).
\]



\section{Criteriu de măsurabilitate pentru vectori aleatori}

\begin{proposition}
O aplicație $X = (X_1, \dots, X_n)$ este vector aleator dacă și numai dacă
fiecare componentă $X_i$ este variabilă aleatoare.

Echivalent:



\[
X \text{ este măsurabil } \Longleftrightarrow
X_i \text{ este măsurabil pentru toate } i = 1,\dots,n.
\]


\end{proposition}

\section{Interpretare}

\begin{itemize}
    \item O variabilă aleatoare este o funcție reală măsurabilă.
    \item Un vector aleator este o funcție cu valori în $\mathbb{R}^n$,
          măsurabilă în sensul produsului.
    \item Măsurabilitatea se verifică doar pe o familie generatoare,
          ceea ce simplifică mult demonstrațiile.
\end{itemize}

\section{Proprietăți suplimentare ale variabilelor aleatoare}

\begin{proposition}
Fie $(\Omega, \mathcal{F}, \mathbb{P})$ un spațiu de probabilitate,
$X : \Omega \to \mathbb{R}$ o variabilă aleatoare
și $g : \mathbb{R} \to \mathbb{R}$ o funcție măsurabilă față de
$\mathcal{B}(\mathbb{R})$.
Atunci compunerea



\[
g \circ X : \Omega \to \mathbb{R}
\]



este variabilă aleatoare.
\end{proposition}

\begin{proposition}
Fie $X : \Omega \to S$ o aplicație măsurabilă între
$(\Omega, \mathcal{F})$ și un spațiu măsurabil $(S, \mathcal{S})$.
Pentru orice subfamilie $\mathcal{C} \subseteq \mathcal{S}$ avem:



\[
\sigma(X^{-1}(\mathcal{C})) = X^{-1}(\sigma(\mathcal{C})).
\]


\end{proposition}

\begin{corollary}
Pentru a verifica că o aplicație $X : \Omega \to S$ este măsurabilă,
este suficient să verificăm că



\[
X^{-1}(C) \in \mathcal{F}
\quad \text{pentru toate } C \in \mathcal{C},
\]



unde $\mathcal{C}$ este o familie de mulțimi care generează
$\sigma$-algebra $\mathcal{S}$ a spațiului țintă.
\end{corollary}

\subsection{Variabile aleatoare extinse}

\begin{definition}
O aplicație



\[
X : \Omega \to \overline{\mathbb{R}} := \mathbb{R} \cup \{-\infty, +\infty\}
\]



se numește \textbf{variabilă aleatoare extinsă} dacă este măsurabilă
față de $\mathcal{F}$ și $\mathcal{B}(\overline{\mathbb{R}})$,
sigma-algebra boreliană extinsă.
\end{definition}

\begin{proposition}
Fie $X, Y : \Omega \to \mathbb{R}$ variabile aleatoare.
Atunci următoarele aplicații sunt variabile aleatoare:



\[
X + Y, \qquad X \cdot Y, \qquad \frac{X}{Y} \ \text{ pe mulțimea } \{Y \neq 0\}.
\]


\end{proposition}

\begin{proposition}
Fie $(X_n)_{n \ge 1}$ un șir de variabile aleatoare extinse.
Atunci următoarele aplicații sunt variabile aleatoare extinse:



\[
\inf_{n \ge 1} X_n, \qquad
\sup_{n \ge 1} X_n, \qquad
\liminf_{n \to \infty} X_n, \qquad
\limsup_{n \to \infty} X_n, \qquad
\lim_{n \to \infty} X_n \ \text{ (când există)}.
\]


\end{proposition}

\end{document}
